% Source: Supplement: residual network and Dinic blocking flow; https://cp-algorithms.com/graph/dinic.html
% Review and validation: docs/validation.md.
\subsection{Dinic 最大流与最小割}
\label{algo:dinic}
\AlgorithmInfo
    {求有向容量网络的最大流，并在残量图上恢复最小割。}
    {一般图时间 $O(n^2m)$，空间 $O(n+m)$；二分匹配的单位网络可达 $O(m\sqrt n)$。}
    {容量为非负 64 位整数；本次增广上限默认 $4\times10^{18}$；源汇不同。}

\subsubsection*{Dinic 的分层与阻塞流}
每轮只沿正残量边 BFS，得到层数 $dep$，再沿 $dep_v=dep_u+1$ 的边推流。阻塞流使当前层次图的每条源汇路径至少有一条饱和边；下一轮汇点的最短层数严格增加，所以轮数小于 $n$。

当前弧指针 \InlineCode{cur[u]} 记住未排除的出边，避免在同一轮中反复扫描无效边。

递归函数 \InlineCode{send} 到达汇点时返回允许流量，回溯时修改正反残量；一条边成功推流后保留当前弧位置，只有走不通时才继续下一条边。一次阻塞流的保守上界为 $O(nm)$，递归深度不超过层数 $O(n)$。残量边与最小割原理见第~\ref{math:flow-basics}~节。

\subsubsection*{使用教程}
\InlineCode{init(n)} 后用 \InlineCode{add(u,v,c)} 加有向容量边，返回正向边编号。\InlineCode{flow(s,t,limit)} 从源 $s$ 向汇 $t$ 最多再送 \InlineCode{limit} 单位，返回\textbf{本次额外增广量}并修改残量容量；省略上限时最多增广 \InlineCode{INF}。同一源汇重复调用从当前残量图继续；若要从零计算，需重新 \InlineCode{init} 并加边。

初始反向容量为零时，原边净流量为 \InlineCode{e[id^1].cap}；\InlineCode{reachable(s)} 返回正残量可达标记。仅在真正跑满最大流后，可据此解释最小割；达到人为限制就停止时不保证是最小割。
\begin{lstlisting}[style=mtycpp,numbers=none]
Dinic g;
g.init(4);
int id=g.add(1,2,3);
g.add(1,3,2),g.add(2,4,2),g.add(3,4,3),g.add(2,3,1);
long long ans=g.flow(1,4);
auto side=g.reachable(1);
long long used=g.e[id^1].cap;
\end{lstlisting}

\needspace{7\baselineskip}
\subsubsection*{模板代码}
\CodeFile{code/graphs/dinic.cpp}

若题目最大流可能超过默认上限，需按数据范围调整上限并保证累加不溢出。无向容量边通常加两条原始有向边，每条都带自己的零容量反向残量边；不要把反向残量边直接初始化为原图容量后仍按普通有向边的流量含义读取。
